% Standalone source-record documentation; no research attempts.
\documentclass[11pt]{article}
\usepackage[margin=25mm]{geometry}
\usepackage{fontspec,unicode-math}
\setmainfont{Latin Modern Roman}
\setmathfont{Latin Modern Math}
\usepackage{amsmath,amssymb,mathtools,xurl,hyperref}
\hypersetup{colorlinks=true,urlcolor=blue,linkcolor=blue}
\setcounter{secnumdepth}{0}
\setlength{\parindent}{0pt}
\setlength{\parskip}{5pt}
\setlength{\emergencystretch}{3em}
\newcommand{\sref}[2]{\hyperlink{source:#1:#2}{[S#2]}}
\newcommand{\cataloguescope}[1]{\paragraph{Source record.}#1}
\newcommand{\floor}[1]{\left\lfloor#1\right\rfloor}
\begin{document}
\section{5.8 (Schleimer) — Why SnapPy works in practice}\label{statement}

\subsection{Definitions and mathematical statement}

\paragraph{Shared notation.}Standard number systems are $\mathbb N,\mathbb Z,\mathbb Q,\mathbb R,\mathbb C$. Specialized notation and index ranges follow the cited source and the definitions below.


A triangulation describes a space by simplices and face identifications; an ideal triangulation places designated vertices at ends or cusps. In hyperbolic $3$-manifold computations, a tetrahedron can be described by a complex shape parameter, and gluing equations impose compatibility around edges and peripheral curves. Numerical procedures approximate solutions to these equations; a certificate requires a specified exact or interval-based verification condition. SnapPy is software working with such triangulated $3$-manifolds and related topological/geometric data.
\paragraph{Record scope and qualifications.}
The catalogue title is retained as metadata. Named certificates, contracts, bounds or frameworks in that title are not asserted as results of the original source extract.
The 2015 list attributes this agenda to Schleimer. It does not specify a SnapPy algorithm, input distribution, success metric or theorem. The context definitions describe terminology from the official documentation; no guarantee for all SnapPy outputs or explanation of its observed performance is claimed.
\paragraph{Statement recorded in the supplied source.}
Give a rigorous explanation for why SnapPy works so well in practice. \sref{statement}{2}
\paragraph{Source status.}
The supplied catalogue labels this source record open; that metadata was checked against the supplied snapshot on 2026-10-11. The label is preserved for provenance and is not an independent verification of present open status or mathematical correctness.
\subsection{Short English statement}
Explain, in a rigorous and explicitly chosen computational model, the practical effectiveness of SnapPy.
\subsection{Sources}
\begin{itemize}
\item[\hypertarget{source:statement:1}{S1}] Ulam AI and the UnsolvedMath Contributors, UnsolvedMath: A Curated Collection of Open Mathematics Problems, catalogue entry. Catalogue curation and metadata: CC BY 4.0. The source extract is retained; definitions, English summaries and qualifications are editorial additions.
\url{https://huggingface.co/datasets/ulamai/UnsolvedMath}.
\item[\hypertarget{source:statement:2}{S2}] AMR Open Problem Lists, source linked by the catalogue.
\url{https://arxiv.org/pdf/1512.04620#page=4}.
\item[\hypertarget{source:statement:3}{S3}] SnapPy official documentation, Triangulation: triangulations and gluing equations. Consulted for terminology, not a universal performance theorem.
\url{https://snappy.computop.org/triangulation.html}.
\end{itemize}
\paragraph{Preparation and review.}
The original source and UnsolvedMath contributors retain attribution for the source extract. Codex (GPT-6; exact serving revision not exposed) assisted with definitions, English summaries, scope notes and formatting. No new human mathematical authorship or independently verified solution is asserted. The supplied catalogue credits: Schleimer. Catalogue identity checked against the complete supplied registry and current repository. Distinct upstream records, including repeated questions and broad or damaged extracts, are retained for complete coverage of the source CS category. A repeated record is not a new independent conjecture. Complete mathematical-equivalence and current-literature review has not been performed. Preparation and final review follow contribution prompts C and D, with source gaps disclosed. Structural and rendering checks do not establish a new theorem or open status.
\label{end:statement}
\end{document}
